\renewcommand{\mydate}{2025年05月27日\ 二十四次课}


%\setcounter{framenumber}{30}
\makeatletter
\ifodd\value{page}
\else
\begin{frame}[plain]
\end{frame}
\fi
\makeatother


\section{内积空间及变换}

\title{第八章 \quad 内积空间及变换}
\author{}
\date{}

\frame{\titlepage}

\subsection{欧氏空间的定义}

\subsubsection{距离与夹角}
\begin{frame}\frametitle{距离与夹角}
三维空间中有距离,夹角等概念 $\xRightarrow{\text{坐标系}}$ 三维数组向量空间$\bR$上向量的模长和夹角.
\pp

比较容易地推广到$\bR^n$: $|(x_1,\cdots,x_n)|:=\sqrt{x_1^2+\cdots+x_n^2}$.
\pp

问题: 如何推广到抽象的线性空间上呢?
\pp
\begin{table}[]
\begin{tabular}{||r||c||l||}\hline\hline
线性空间 & $\xleftrightarrow[1:1]{\text{基} \alpha_1,\cdots,\alpha_n}$ & $\bR^n$ \\ \hline\hline
向量$v$ & $\longleftrightarrow$ & 坐标 \\ \hline
线性变换$\sA$ & $\longleftrightarrow$ & A \\ \hline
?? & $\longleftrightarrow$ & 长度,夹角 \\ \hline
$\sA$(保持长度) & $\longleftrightarrow$ & ?? \\\hline
$\sA$(??) & $\longleftrightarrow$ & 转置矩阵 \\ \hline
$\sA$(??) & $\longleftrightarrow$ & 对称矩阵 \\ \hline\hline
\end{tabular}
\end{table}
\end{frame}






\begin{frame}\frametitle{距离与夹角}
回顾$\bR^3$上的模长与夹角公式:
\pp
\begin{equation*}
\begin{split}
|a| & = \sqrt{(a,a)}\\
\pp
\theta & = \arccos\frac{(a,b)}{\sqrt{(a,a)(b,b)}}
\end{split}
\end{equation*}
\pp
容易看出模长和夹角可以由内积完全确定.
\pp
反之,内积可由模长唯一确定
\[(a,b)=\frac{|a+b|^2-|a|^2-|b|^2}{2}.\]
\pp
总结 $\bR^3$上内积有如下基本性质:
\begin{enumerate}
\item 对称性
\item 线性性
\item 正定性
\end{enumerate}
\pp
\yang{为了在一般的$\bR$-线性空间上定义度量,我们需要引入一般空间上的内积.}
\end{frame}








\subsubsection{欧氏空间的定义}
\begin{frame}\frametitle{欧氏空间的定义}
\begin{definition}[内积,欧氏空间]
\pp
设$V$为有限维$\bR$-线性空间.
\pp
若存在一个映射$(-,-)\colon V\times V \rightarrow \bR$满足:
\pp
\begin{enumerate}
\item 对称性: $(a,b)=(b,a)$;
\pp
\item 双线性: $(\lambda a,b)=\lambda(a,b)$, $(a+b,c)=(a,c)+(b,c)$;
\pp
\item 正定性: $(a,a)\geq0$,且等号成立当且仅当$a=0$,
\end{enumerate}
\pp
则称$(a,b)$为$a$和$b$的\blue{内积}.
\pp 带内积的$\bR$-线性空间称为\blue{欧氏空间}(或,\blue{欧几里得空间}).
\end{definition}
\pp
\begin{center}
欧式空间 = $\bR$-线性空间 + 内积;\\
$\bR$-线性空间 = 集合 + 加法 + ($\bR$)数乘\\
\end{center}
\pp
注: 内积$(-,-)$对第二个分量也是线性的.
\pp
\begin{proposition}[内积的基本性质]
\begin{enumerate}
\item $\left(\sum_{i=1}^{k}\lambda_ia_i,\sum_{j=1}^{m}\mu_jb_j\right) = \sum_{i=1}^{k}\sum_{j=1}^{m} \lambda_i\mu_j(a_i,b_j)$;
\pp
\item $(a,0)=0$;
\pp
\item 内积结构由一组基之间的内积完全确定.
\end{enumerate}
\end{proposition}
\end{frame}





\subsubsection{模长与夹角}
\begin{frame}\frametitle{模长与夹角}
\pp

为了定义夹角,我们需要引入:

\begin{theorem}[Cauchy-Schwarz 不等式]
\pp
设$V$为欧式空间. 对于任意$a,b\in V$,我们有
\[|(a,b)|\leq \sqrt{(a,a)\cdot(b,b)}.\]
\end{theorem}
\pp
\yang{证明思路:考虑二次多项式$(xa+b,xa+b)$的判别式.}

\pp
\yang{可以定义夹角了!(不确定原理.)}
\pp
\begin{definition}[长度,夹角]
设$V$为欧氏空间. 对于任意向量 $a,b\in V$,
\begin{enumerate}
\item 称$|a|:=\sqrt{(a,a)}$为$a$的\blue{长度}(或,\blue{模长});
\pp
\item 称$|a-b|$为$a$和$b$之间的\blue{距离}.记为\blue{$d(a,b)$}.
\pp
\item 若$a,b$不为零,定义$a$和$b$之间的\blue{夹角}为 $\theta = \arccos\frac{(a,b)}{|a|\cdot|b|}$.
\pp
特别地,当$(a,b)=0$时,称$a$与$b$\blue{正交}或\blue{垂直},记为 \blue{$a\perp b$}.
\pp
\item 称$a$为\blue{单位向量},若$|a|=1$.
\pp
若$b\neq0$,称$\frac{1}{|b|}b$为$b$的\blue{单位化}.
\end{enumerate}
\end{definition}
\end{frame}





\subsubsection{距离基本性质}
\begin{frame}\frametitle{距离基本性质}
\pp
\begin{proposition}
\begin{enumerate}
\item 对称性: $d(a,b)=d(b,a)$;
\pp
\item 正定性: $d(a,b)\geq0$, 等号成立当且仅当$a=b$;
\pp
\item 三角不等式: $d(a,c)\leq d(a,b)+ d(b,c)$;
\pp
\item 勾股定理: $a\perp b$ 当且仅当 $|a|^2+|b|^2=|a+b|^2$;
\pp
\item 菱形定理: 若$|a|=|b|$,则$(a+b)\perp(a-b)$.
\end{enumerate}
\end{proposition}
\end{frame}





\begin{frame}\frametitle{内积的基本例子}
\pp
\begin{example}
设$V=\bR^n$,对于任意两向量 $a=(a_1,\cdots,a_n)$和 $b=(b_1,\cdots,b_n)$, 定义
\[(a,b)=a_1b_1+\cdots+a_nb_n.\]
\pp
则$(-,-)$为$V$上的一个内积.这一内积给出通常的模长和夹角.
\pp
此时Cauchy-Schwarz公式为
\[\left(\sum_{i=1}^n a_ib_i\right)^2 \leq \left(\sum_{i=1}^n a_i^2\right)\cdot\left(\sum_{i=1}^n b_i^2\right).\]
\end{example}
\pp
\begin{example}
类似的$(a,b)=a_1b_1+2a_2b_2+\cdots+na_nb_n$也可以定义$V=\bR^n$上的一个内积.
\end{example}
\end{frame}






\begin{frame}\frametitle{内积的基本例子}

\yang{构造一个内积使得给定的一组基均为单位向量且两两垂直.}
\pp
\begin{example}
设$\alpha_1,\cdots,\alpha_n$为$n$维$\bR$-线性空间$V$上的一组基.
\pp
对于任意向量$\alpha,\beta\in V$,定义
\[(\alpha,\beta)=a_1b_1+a_2b_2+\cdots+a_nb_n,\]
其中$(a_1,a_2,\cdots,a_n)$和$(b_1,b_2,\cdots,b_n)$分别为$\alpha$和$\beta$在基$\alpha_1,\cdots,\alpha_n$下的坐标.
\pp
则$(-,-)$为$V$上的内积.在此内积下$\alpha_1,\cdots,\alpha_n$均为单位向量且两两垂直.
\pp
即,
\[(\alpha_i,\alpha_j)=\delta_{ij}=\left\{\begin{array}{ll}
1,& i=j;\\
0,& i\neq j.\\
\end{array}\right.\]
\end{example}
\pp
\yang{任意有限维实线性空间一定存在内积. 注:任意内积均为这一形式!}
\end{frame}




\subsubsection{Hilbert空间*}
\begin{frame}\frametitle{Hilbert空间*}

在无穷维空间上也可类似的定义内积.
\pp
\begin{example}[无限维空间上的内积]
对应任意$[a,b]$上的连续函数$f,g\in C[a,b]$,定义
\[(f,g):=\int_a^b f(x)g(x)\mathrm{d}x.\]
\pp
类似地,定义连续函数的模长,夹角和正交等概念.此时Cauchy-Schwarz公式为
\[\left(\int_a^b f(x)g(x)\mathrm{d}x\right)^2\leq\left(\int_a^b f(x)^2\mathrm{d}x\right)\cdot\left(\int_a^b g(x)^2\mathrm{d}x\right).\]
\pp
当$a=-\pi$,$b=\pi$时,三角函数$1,\cos x,\sin x,\cos2x,\sin2x,\cdots$在这一内积下两两正交.
\end{example}
\pp
我们称带内积\blue{完备的}$\bR$-线性空间为\blue{Hilbert空间}.

\end{frame}



\subsection{度量矩阵(内积的矩阵表示)}

\subsubsection{内积的矩阵表示}
\begin{frame}\frametitle{内积的矩阵表示}
\pp
设$\alpha_1,\cdots,\alpha_n$为欧氏空间$V$的一组基.
\pp
任取$\alpha = a_1\alpha_1+\cdots+a_n\alpha_n$, $\beta = b_1\alpha_1+\cdots+b_n\alpha_n\in V$,则
\begin{equation}\label{equ:innerprod}
(\alpha,\beta)=\sum_{i=1}^n\sum_{j=1}^n a_ib_j(\alpha_i,\alpha_j).
\end{equation}
\pp
因此内积$(-,-)$由值$g_{ij}:=(\alpha_i,\alpha_j)$($1\leq ij\leq n$)唯一确定.
\pp
记
\[G:=(g_{ij})_{n\times n}.\]
称$G$为内积$(-,-)$在基$\alpha_1,\cdots,\alpha_n$下的\blue{度量矩阵}.
\pp
记
\[x=\begin{pmatrix}a_1\\\vdots\\a_n\end{pmatrix},y=\begin{pmatrix}b_1\\\vdots\\b_n\end{pmatrix}\]
则\eqref{equ:innerprod}可简写为
\[\blue{(\alpha,\beta)=x^TGy.}\]
\end{frame}


\renewcommand{\mydate}{2025年05月29日\ 二十五次课}



\begin{frame}\frametitle{}
通过度量矩阵我们有如下映射
\begin{equation*}
\xymatrix@C=2cm@R=0cm{
\text{$V$上的内积}
\ar[r]^-{\text{基}\alpha_1,\cdots,\alpha_n}	& \bR\text{上的$n$阶矩阵} \\
(-,-) \ar@{|->}[r] & G=\Big((\alpha_i,\alpha_j)\Big)_{n\times n} \\
}
\end{equation*}
\pp

\yang{问题: 哪些矩阵落在这个映射的像集里面?即,哪些矩阵能够成为某个内积的度量矩阵?}
\pp

\begin{proposition}[度量矩阵的基本性质]
\begin{enumerate}
\item 设$G$为$V$上某内积$(-,-)$在基$\alpha_1,\cdots,\alpha_n$下的度量矩阵.
\pp
则
\begin{itemize}
\item $G$为实对称矩阵;
\pp
\item 对任意$x\in\bR^n$,都有$x^TGx\geq0$,且等号成立当且仅当$x=0$.
\end{itemize}
\pp
称满足如上性质的矩阵为\blue{正定矩阵}.
\pp
因此内积的度量矩阵为正定矩阵.
\pp
\item 反之,对于任意给定的正定矩阵$G$, 通过
\[(\alpha,\beta):=x^TGy\]
可以构造出$V$上的一个内积,其中$x,y$为$\alpha$和$\beta$的坐标.
\end{enumerate}
\end{proposition}
\end{frame}






\begin{frame}\frametitle{}
\yang{问题:
\begin{enumerate}
\item 不同基下度量矩阵之间的关系?
\pp
\item 度量矩阵的最简形式,在哪组基下度量矩阵为最简单形式?
\end{enumerate}
}
\pp
\begin{proposition}
设$P$为欧氏空间的两组基$\alpha_1,\cdots,\alpha_n$和$\eta_1,\cdots,\eta_n$为之间的过渡矩阵
\[(\eta_1,\cdots,\eta_n) = (\alpha_1,\cdots,\alpha_n)P.\]
\pp
设内积在$\alpha_1,\cdots,\alpha_n$和$\eta_1,\cdots,\eta_n$下的度量矩阵为$G$和$\overline{G}$. 即,
\[G=\Big((\alpha_i,\alpha_j)\Big)_{n\times n},\quad \overline{G}=\Big((\eta_i,\eta_j)\Big)_{n\times n}.\]
\pp
则
\[\overline{G}=P^TGP.\]
\end{proposition}
\pp
\yang{证明思路: $\overline{G}_{(i,j)}=(\eta_i,\eta_j)=\left(\sum\limits_{s=1}^{n}p_{si}\alpha_s,\sum\limits_{t=1}^{n}p_{tj}\alpha_t\right)=\sum\limits_{s=1}^n\sum\limits_{t=1}^n p_{si}G_{(s,t)} p_{tj} = (P^TGP)_{(i,j)}$.}
\end{frame}





\subsubsection{相合关系及其基本性质}
\begin{frame}\frametitle{相合关系及其基本性质}
\begin{definition}[相合]
称两个矩阵$G$和$\overline{G}$ \blue{相合},若存在可逆阵$P$使得
\[\overline{G}=P^TGP.\]
\end{definition}
\pp
\begin{proposition}
\begin{enumerate}
\item 内积在不同基下的度量矩阵相合;
\pp
\item 相合为等价关系.
\end{enumerate}
\end{proposition}
\pp
实对称矩阵的相合分类,以及相合标准形(第八章).

\end{frame}




\subsection{标准正交基}


\subsubsection{标准正交基}
\begin{frame}\frametitle{标准正交基}
\yang{度量矩阵的最简形式?为了回答这一问题,我们需要引入标准正交基.}

\pp
\begin{definition}[标准正交基]
设$V$为$n$维欧氏空间.
\begin{enumerate}
\item 称由一组两两正交的非零向量为\blue{正交向量组};
\pp
\item 称由正交向量组构成的基为\blue{正交基};
\pp
\item 称由单位向量组成的正交基为\blue{标准正交基}.
\end{enumerate}
\end{definition}
\pp
\begin{example}设$(-,-)$为$\bR^n$上的一内积.
\pp
\begin{enumerate}
\item 若$((x_1,\cdots,x_n),(y_1,\cdots,y_n))=\sum_{i=1}^nx_iy_i$.
\pp
则自然基
\[e_1,\cdots,e_n\]
为一组标准正交基.
\pp
\item 若$((x_1,\cdots,x_n),(y_1,\cdots,y_n))=\sum_{i=1}^nix_iy_i$.
\pp
则
\[e_1,\frac{e_2}{\sqrt{2}}\cdots,\frac{e_n}{\sqrt{n}}\]
为一组标准正交基.
\end{enumerate}
\end{example}
\end{frame}









% ========== 2023年06月06号 第27次课 标准正交基,正交变换 ==========
%\frame{\titlepage}






\begin{frame}\frametitle{标准正交基}
\yang{度量矩阵的最简形式?为了回答这一问题,我们需要引入标准正交基.}

\begin{definition}[标准正交基]
设$V$为$n$维欧氏空间.
\begin{enumerate}
\item 称由一组两两正交的非零向量为\blue{正交向量组};
\item 称由正交向量组构成的基为\blue{正交基};
\item 称由单位向量组成的正交基为\blue{标准正交基}.
\end{enumerate}
\end{definition}

\begin{proposition}
正交向量组线性无关.
\end{proposition}
\pp

\begin{theorem}[标准正交基使度量矩阵最简]
设$V$为$n$维欧氏空间. 设$G$为内积在基$\alpha_1,\cdots,\alpha_n$下的度量矩阵.
\pp
则
\begin{center}
$\alpha_1,\cdots,\alpha_n$为标准正交基 \quad $\Leftrightarrow$ \quad $G=I_n$.
\end{center}
\end{theorem}
\pp
下面将讨论标准正交基的存在性.
\end{frame}





\subsubsection{Schmidt正交化}
\begin{frame}\frametitle{Schmidt正交化}

\begin{theorem}[Schmidt正交化]
给定欧氏空间的任意一组基$\alpha_1,\cdots,\alpha_n$,则存在一组标准正交基$e_1,\cdots,e_n$使得(对所有的$i=1,2,\cdots,n$)
\[\langle\alpha_1,\cdots,\alpha_i\rangle = \langle e_1,\cdots,e_i\rangle.\footnote{\text{也等价于 $(\alpha_1,\cdots,\alpha_i) = (e_1,\cdots,e_i)R$, 其中$R$为可逆上三角矩阵.}}\]

\end{theorem}
\pp
\yang{几何解释?}
\yang{证明思路: 递归地定义 $\beta_k=\alpha_k - \sum_{i=1}^{k-1}(\alpha_k,e_i)e_i\neq0$ 以及 $e_k=\frac{\beta_k}{|\beta_k|}$.}
\pp

\begin{corollary}[实可逆阵的QR分解]
对任意可逆实矩阵 $A$, 存在方阵$Q$和上三角矩阵$R$使得$A=QR$, 其中$Q$的列向量构成$\mathbb R^n$的标准正交基.
\end{corollary}
\end{frame}



\renewcommand{\mydate}{2025年06月03日\ 二十六次课}


\begin{frame}\frametitle{Schmidt正交化}
\begin{example}
将
$\alpha_1=\begin{pmatrix}1\\1\\0\\0\end{pmatrix}$,
$\alpha_2=\begin{pmatrix}1\\0\\1\\0\end{pmatrix}$,
$\alpha_3=\begin{pmatrix}-1\\0\\0\\1\end{pmatrix}$,
$\alpha_4=\begin{pmatrix}1\\-1\\-1\\0\end{pmatrix}$标准正交化.
\end{example}
\pp
\begin{corollary}
对任意正定矩阵$A$,存在可逆实矩阵$P$使得$A=P^TP$.
\end{corollary}
\end{frame}



\subsection{正交变换}

\subsubsection{正交变换}
\begin{frame}\frametitle{正交变换}
\begin{table}[]
\begin{tabular}{||r||c||l||}\hline\hline
$n$维欧氏空间 & $\xleftrightarrow[1:1]{\text{标准正交基基} \alpha_1,\cdots,\alpha_n}$ & $\bR^n$(带标准内积) \\ \hline\hline
线性变换$\sA$ & $\xleftrightarrow[1:1]{\sA(\alpha_1,\cdots,\alpha_n)=(\alpha_1,\cdots,\alpha_n)A}$ & 矩阵A \\ \hline
$\sA$(保持内积) & $\longleftrightarrow$ & ?? \\\hline
\end{tabular}
\end{table}

\pp
\begin{example}空间(平面)的旋转和镜面反射等.
\end{example}
\pp

\begin{definition}[正交变换]
设$\sA$为欧氏空间$V$上的线性变换.若$\sA$ \CJKunderdot{保持内积},即(对任意$a,b\in V$,
\[(\sA a,\sA b)=(a,b),\]
则称$\sA$为\blue{正交变换}.
\end{definition}
\end{frame}





\subsubsection{正交变换的等价刻画}
\begin{frame}\frametitle{正交变换的等价刻画}
\pp
\begin{theorem}
设$\sA$为欧氏空间$V$上的线性变换.则以下几条等价
\begin{enumerate}
\item $\sA$ 正交;
\pp
\item $\sA$ 保持向量长度;
\pp
\item $\sA$ 将标准正交基变为标准正交基.
\end{enumerate}
\end{theorem}
\pp
\yang{\tiny 证明思路:(1)$\Leftrightarrow$(2) \& (1)$\Leftrightarrow$(3)

\pp
(1)$\Rightarrow$(2): $|\sA(\alpha)|=\sqrt{(\sA(\alpha),\sA(\alpha))}=\sqrt{(\alpha,\alpha)}=|\alpha|;$
\pp

(2)$\Rightarrow$(1): $(\sA(\alpha),\sA(\beta)) = \frac{|\sA(\alpha+\beta)|^2-|\sA(\alpha)|^2-|\sA(\beta)|^2}{2}$
\[= \frac{|\alpha+\beta|^2-|\alpha|^2-|\beta|^2}{2}
=(\alpha,\beta)\]
\pp
(1)$\Rightarrow$(3):
$e_1,\cdots,e_n$ = 标准正交基 $\Rightarrow$ $(\sA(e_i),sA(e_j))=(e_i,e_j)=\delta_{ij}$
\begin{center}
$\Rightarrow$ $\sA(e_1),\cdots,\sA(e_n)$ = 标准正交基.
\end{center}
\pp
(3)$\Rightarrow$(1):设$\sA$将标准正交基$e_1,\cdots,e_n$映为另一组标准正交基$\sA(e_1),\cdots,\sA(e_n)$. 任取
\[\alpha = \sum_{i=1}^n a_ie_i\quad \text{和}\quad \beta = \sum_{i=1}^{n}b_ie_i.\]
则
\begin{equation*}
\begin{split}
(\sA(\alpha),\sA(\beta)) & = \left(\sum_{i=1}^n a_i\sA(e_i),\sum_{i=1}^{n}b_j\sA(e_j)\right) = \sum_{i=1}^n\sum_{j=1}^{n}a_ib_j\left(\sA(e_i),\sA(e_j)\right)\\
&= \sum_{i=1}^n\sum_{j=1}^{n}a_ib_j\delta_{ij}= \sum_{i=1}^n\sum_{j=1}^{n}a_ib_j(e_i,e_j)= (\alpha,\beta)\\
\end{split}
\end{equation*}
}
\end{frame}





\subsubsection{正交变换基本性质}
\begin{frame}\frametitle{正交变换基本性质}
\pp
\begin{theorem}[全体正交变换在复合运算下构成群]
\pp
设$V$为欧氏空间.
则
\begin{enumerate}
\item 单位变换为正交变换;
\pp
\item 正交变换的复合仍然为正交变换;
\pp
\item 正交变换可逆且其逆也为正交变换.
\end{enumerate}
\end{theorem}
\pp
\yang{证明思路: 保持长度.}

\pp
\yang{注: 正交变换群描述了欧氏空间的对称性.}
\end{frame}





\subsubsection{正交变换在标准正交基下的矩阵}
\begin{frame}\frametitle{正交变换在标准正交基下的矩阵}
设正交变换$\sA$在标准正交基$e_1,\cdots,e_n$下的矩阵为$A$.

\yang{这个矩阵$A$满足什么特别的性质?}
\pp
\begin{equation*}
\begin{split}
\sA\text{正交}
& \Leftrightarrow \sA e_1,\cdots,\sA e_n \text{为标准正交基}\\
\pp
& \Leftrightarrow (\sA e_i,\sA e_j)=\delta_{ij} (\text{对任意的}1\leq i,j\leq n)\\
\pp
& \Leftrightarrow \left(\sum_{\ell=1}^{n} a_{\ell i}e_\ell,\sum_{k=1}^na_{kj} e_k\right)=\delta_{ij} (\text{对任意的}1\leq i,j\leq n)\\
\pp
& \Leftrightarrow \sum_{k=1}^{n} a_{ki}a_{kj}=\delta_{ij} (\text{对任意的}1\leq i,j\leq n)\\
\pp
& \Leftrightarrow A^TA=I_n\\
\end{split}
\end{equation*}
\pp
\begin{definition}[正交矩阵]
若$n$阶实方阵$A$满足$A^TA=I_n$(或$A^{-1}=A^T$),则称$A$为\blue{正交矩阵}.
\end{definition}
\pp
注:正交矩阵的行向量(或列向量)构成$\bR^n$的标准正交基(在标准内积下).
\end{frame}









% ========== 2023年06月08号 第28次课 正交变换与对称变换 ==========
%\frame{\titlepage}






\begin{frame}\frametitle{正交变换在标准正交基下的矩阵}
\begin{theorem}
设$\sA$为欧氏空间$V$上的线性变换,则$\sA$为正交变换当且仅当$\sA$在标准正交基下的矩阵$A$为正交矩阵.
\end{theorem}
\pp

\begin{center}
正交变换$\sA$ $\xleftrightarrow[1:1]{\text{标准正交基}}$ 正交矩阵$A$
\end{center}
\pp

\begin{definition}[第一类变换,第二类变换]
设$A$为正交矩阵.则 $A^TA=I$.
\pp
因此$\det(A)=\pm1$.
\pp
\begin{enumerate}
\item 若$\det(A)=1$,则称$\sA$为\blue{第一类变换};
\pp
\item 若$\det(A)=-1$,则称$\sA$为\blue{第二类变换}.
\end{enumerate}
\end{definition}
\pp
\begin{example}
三维空间或在二维空间的旋转为第一类变换,而镜面反射为第二类变换.
\end{example}
\end{frame}






\begin{frame}\frametitle{}
\begin{proposition}
设$\sA$为欧氏空间$V$上的正交变换.
\pp
\begin{enumerate}
\item 则$\sA$的特征值模长都为$1$.特别地,实特征值只可能为$\pm1$.
\pp
\item 若$V$的维数为奇数且$\sA$为第一类正交变换,则$1$为$\sA$的特征值.
\end{enumerate}
\end{proposition}
\pp
\yang{证明思路: (1). 设$A$为$\sA$在某组标准正交基下的矩阵.设$A(\xi)=\lambda\xi$,则
\[\overline{\xi}^T\xi = \overline{\xi}^T\overline{A}^TA\xi=\overline{\lambda}\lambda \overline{\xi}^T\xi.\]
因此 $|\lambda|=1$.
\pp

(2).多项式$\varphi_{\sA}(\lambda)$为实系数,其复根共轭成对出现.而实根仅取$\pm1$. 若$1$不为特征值,则$-1$出现奇数次,因此行列式为$-1$.矛盾!
}
\pp
\begin{corollary}
三维空间中的第一类正交变换保持一个向量不变,从而一定为旋转变换.
\end{corollary}
\end{frame}


\renewcommand{\mydate}{2025年06月05日\ 二十七次课}


\subsection{伴随变换}


\subsubsection{转置与伴随变换}
\begin{frame}\frametitle{转置与伴随变换}
\begin{table}[]
\begin{tabular}{||r||c||l||}\hline\hline
$n$维欧氏空间 & $\xleftrightarrow[1:1]{\text{标准正交基基} \alpha_1,\cdots,\alpha_n}$ & $\bR^n$(带标准内积) \\ \hline\hline
线性变换$\sA$ & $\xleftrightarrow[1:1]{\sA(\alpha_1,\cdots,\alpha_n)=(\alpha_1,\cdots,\alpha_n)A}$ & 矩阵$A$ \\ \hline
正交变换$\sA$ & $\xleftrightarrow[1:1]{\sA(\alpha_1,\cdots,\alpha_n)=(\alpha_1,\cdots,\alpha_n)A}$ & 正交矩阵$A$ \\ \hline
?? & $\longleftrightarrow$ & 转置矩阵 $A^T$ \\\hline
\end{tabular}
\end{table}
\pp
\begin{theorem}
设$\sA$为欧氏空间$V$上的线性变换.
设$\sA$在标准正交基$e_1,\cdots,e_n$下的矩阵为$A$.
\pp
则
\begin{enumerate}
\item \blue{存在唯一}的$V$上的线性变换,记为$\sA^*$,满足
\[(\sA \alpha,\beta) =(\alpha,\sA^*\beta), \quad(\forall \alpha,\beta\in V).
\]
\pp
\item 线性变换$\sA^*$在基$e_1,\cdots,e_n$下的矩阵为$A^T$.
\end{enumerate}
\pp
称$\sA^*$为$\sA$的\blue{伴随变换}.
\end{theorem}
\pp
\yang{
证明思路: 存在性: 证明矩阵$A^T$对应的线性变换满足性质. 唯一性: 若$\sB$也满足条件,则仅需证明 $|(\sA^*-\sB)(\beta)|=0$ 对所有的$\beta$成立.
}
\end{frame}





\subsubsection{伴随变换的基本性质}
\begin{frame}\frametitle{伴随变换的基本性质}
\pp
\begin{proposition}
\begin{enumerate}
\item $(\sA^*)^*=\sA$;
\pp
\item $(\sA+\sB)^*=\sA^*+\sB^*$;
\pp
\item $(\lambda \sA)^* = \lambda \sA^*$;
\pp
\item $(\sA\circ\sB)^* = \sB^*\circ\sA^*$;
\end{enumerate}
\end{proposition}
\pp
\begin{proposition}
设$\sA$为欧氏空间上的线性变换.则
\begin{center}
$\sA$ 正交 \quad $\Leftrightarrow$ \quad $\sA^*\sA=\varepsilon$ \quad $\Leftrightarrow$ \quad $\sA^*$ 正交.
\end{center}
\end{proposition}
\end{frame}





\subsubsection{对称变换与对称矩阵}
\begin{frame}\frametitle{对称变换与对称矩阵}
\pp
\begin{definition}[对称变换,自伴随变换]
设$\sA$为欧氏空间$V$上的线性变换.
\pp
若$\sA=\sA^*$,则称$\sA$为$V$上的\blue{对称变换}(或\blue{自伴随变换}).
\end{definition}
\pp
注: $\sA$对称 $\Leftrightarrow$ $(\sA a,b)=(a,\sA b)$ $\forall a,b\in V$.

\pp
\begin{theorem}
设$A$为某欧氏空间上的线性变换$\sA$在某组标准正交基下的矩阵.则
\begin{center}
$\sA$ 为对称变换 $\Leftrightarrow$ $A$为实对称矩阵.
\end{center}
\end{theorem}
\pp
\yang{证明: $\sA$对称 $\Leftrightarrow$ $\sA^*=\sA$ $\Leftrightarrow$ $A^T=A$ $\Leftrightarrow$ $A$对称.}

\pp
\begin{theorem}
对称变换的不同特征值对应的特征向量正交.
\end{theorem}
\pp
\yang{证明思路: $(\sA(\xi_1),\xi_2) = (\xi_1,\sA(\xi_2))$ $\Rightarrow$ $(\lambda_1-\lambda_2)\cdot(\xi_1,\xi_2)=0$.}

\pp
\begin{corollary}
实对称矩阵$A$的属于不同特征值的特征向量正交.
\end{corollary}
\end{frame}


\subsection{实对称矩阵的对角化}


\subsubsection{实对称矩阵的对角化}
\begin{frame}\frametitle{实对称矩阵的对角化}
\yang{下面将证明实对称矩阵总是可对角化的.}
\pp
\begin{proposition}
实对称矩阵的特征值都为实数.
\end{proposition}
\pp
\yang{证明思路: 设$A\xi=\lambda\xi$($\xi\neq0$). 由于$A$为实矩阵, $A\overline{\xi}=\overline{A\xi}=\overline{\lambda}\cdot\overline\xi$.
由于$A$对称,我们有$(A\overline\xi)^T\xi = \overline\xi^T(A\xi)$. 因此 $\overline{\lambda}=\lambda$.
}
\pp

\begin{theorem}
任意$n$阶实对称矩阵$A$,存在$n$阶正交矩阵$T$使得$T^{-1}AT$为对角矩阵.
\end{theorem}

\pp
\yang{证明思路:将一个单位特征向量扩充为一组标准正交基,并得正交阵$T_n$.则 $T_n^{-1}AT_n=\begin{pmatrix}\lambda_1&\\&A_{n-1}\\\end{pmatrix}$,其中$A_{n-1}$仍然为实对称.归纳即可得证.}
\pp
\begin{example}
设$A=\begin{pmatrix}1&2&2\\2&1&2\\2&2&1\end{pmatrix}$.求正交矩阵$T$使得$T^{-1}AT$为对角矩阵.
\end{example}
\pp
\yang{提示:$P_A(\lambda)=(\lambda-5)(\lambda+1)^2$.}
\end{frame}





\renewcommand{\mydate}{2025年06月10日\ 二十八次课}


\subsection{欧氏空间的子空间}


\subsubsection{欧氏空间的子空间}
\begin{frame}\frametitle{欧氏空间的子空间}
\pp
\begin{definition}[正交]
设$V_1,V_2$为欧氏空间$V$的两子空间.
\pp
若对任意$a_1\in V_1$,$a_2\in V_2$都有$(a_1,a_2)=0$,则称$V_1$和$V_2$\blue{相互正交},记为\blue{$V_1\perp V_2$}.
\pp
若一个向量$a$满足$\langle a\rangle\perp V_1$,则称$a$与$V_1$\blue{正交},记为\blue{$a\perp V_1$}.
\end{definition}
\pp
\begin{theorem}
\begin{enumerate}
\item 若$V_1\perp V_2$,则$V_1+V_2$为直和;
\pp
\item 若$V_1,V_2,\cdots,V_r$两两正交,则$V_1+V_2+\cdots+V_r$为直和.
\end{enumerate}
\end{theorem}
\pp

\begin{definition}[正交补]
若$V_1\perp V_2$且$V=V_1+V_2$,则称$V_1,V_2$互为\blue{正交补(空间)}.
\end{definition}
\pp

\begin{theorem}
欧氏空间的任意子空间的正交补存在且唯一.
\end{theorem}
\end{frame}



\subsubsection{正交投影}
\begin{frame}\frametitle{正交投影}
\begin{definition}[正交投影]
设\(W\)是欧氏空间\(\mathbb{R}^n\)的子空间, \(\boldsymbol{x} \in \mathbb{R}^n\). 若\(\boldsymbol{y} \in W\)满足\(\boldsymbol{x}-\boldsymbol{y} \perp W\), 则称\(\boldsymbol{y}\)为\(\boldsymbol{x}\)在\(W\)上的\blue{正交投影}, 记为\(\boldsymbol{y} = P\boldsymbol{x}\).
\end{definition}
\pp
\begin{theorem}[存在、唯一]
设\(W\)为欧氏空间\(\mathbb{R}^n\)的子空间. 设\(\boldsymbol{u}_1,\boldsymbol{u}_2,\ldots,\boldsymbol{u}_r\)是$W$的一组标准正交基. \pp 则
\begin{enumerate}
\item 对任意\(\boldsymbol{x} \in \mathbb{R}^n\)在\(W\)上的正交投影存在唯一, 且 \\
$P\boldsymbol{x} = (\boldsymbol{x},\boldsymbol{u}_1)\boldsymbol{u}_1 + (\boldsymbol{x},\boldsymbol{u}_2)\boldsymbol{u}_2 + \ldots + (\boldsymbol{x},\boldsymbol{u}_r)\boldsymbol{u}_r.$ \pp
\item  对任意\(\boldsymbol{y} \in W\),
$|\boldsymbol{x} - \boldsymbol{y}| \geq |\boldsymbol{x} - P\boldsymbol{x}|$.
等号成立当且仅当\(\boldsymbol{y} = P\boldsymbol{x}\).
\end{enumerate}
\end{theorem}
\pp
\begin{example}
\begin{enumerate}
\item 求$(2,1,3)$在$\langle(1,1,0),(1,0,-1)\rangle$上的正交投影. \pp
\item 设\(A \in \mathbb{R}^{m\times n}\), \(\boldsymbol{b} \in \mathbb{R}^m\), 且\(A\)列满秩. 求线性方程组\(A\boldsymbol{x} = \boldsymbol{b}\)的\blue{最小二乘解}, 即求\(\boldsymbol{x} \in \mathbb{R}^n\)使得\(\vert A\boldsymbol{x}-\boldsymbol{b}\vert\)达到最小.
\end{enumerate}
\end{example}
\pp
\yang{考虑$A$的列空间$C(A)$. $\boldsymbol{x}=(A^TA)^{-1}A^T\boldsymbol{b}$.}
\end{frame}



\subsection{酉空间*}

\subsubsection{酉空间*}
\begin{frame}\frametitle{酉空间*}
\begin{equation*}
\xymatrix{
\text{$\bR$-线性空间} \ar@{}[r]|-{+} & \text{内积} \ar@{=>}[r] \ar@{=>}[d]& \text{欧氏空间} \ar@{=>}[d]\ar@/_30pt/@{=>}[dd]\\
&\text{度量矩阵} \ar@{=>}[d] & \text{正交变换} \ar@{=>}[r] & \text{正交矩阵}\\
&\text{正定矩阵} & \text{伴随变换} \ar@{=>}[d] \ar@{=>}[r] & \text{转置}\\
&               & \text{自伴随变换}\ar@{=>}[r] & \text{对称矩阵}\\
}
\end{equation*}
\pp
\begin{equation*}
\xymatrix{
\text{$\bC$-线性空间} \ar@{}[r]|-{+} & \text{内积(??)} \ar@{=>}[r] & \text{??空间} \\
}
\end{equation*}
\end{frame}





\subsubsection{复数模长和复数组向量模长*}
\begin{frame}\frametitle{复数模长和复数组向量模长*}
\pp
\begin{enumerate}
\item $x\in \bR$ $\Rightarrow$ $|x| =\sqrt{x\cdot x}$;
\pp
\item $x=(x_1,\cdots,x_n)\in \bR^n$ $\xRightarrow{\text{标准内积下}}$ $|x|=\sqrt{x_1x_1+\cdots+x_nx_n}$;
\pp
\item $z\in \bC$ $\Rightarrow$ $|z| =\sqrt{\overline{z}\cdot z}$;
\end{enumerate}
\pp
很自然地,我们可以如下定义复数组向量的长度:
\begin{center}
$z=(z_1,\cdots,z_n)\in \bC^n$ $\Rightarrow$ $|z|=\sqrt{\overline{z_1}z_1+\cdots+\overline{z_n}z_n}$;
\end{center}
\pp
对于任意复矩阵$A$,记
\[A^H:= \overline{A}^T\]
称其为$A$的\blue{共轭转置}.
\pp
则复数组(列)向量的长度公式可写为
\[|z|^2 = z^H\cdot z.\]
\pp
注:共轭转置保持加法,但只是在共轭意义下保持数乘.
\end{frame}





\subsubsection{复向量空间上的内积定义*}
\begin{frame}\frametitle{复向量空间上的内积定义*}
\pp
\begin{definition}[酉空间]
设$V$为复线性空间.
\pp
若映射
\[(-,-)\colon V\times V\longrightarrow \bC\]
满足
\pp
\begin{enumerate}
\item 共轭对称性 $(a,b)=\overline{(b,a)}$;
\pp
\item 线性性 $(a,\lambda b + \mu c)=\lambda(a,b)+\mu(a,c)$;
\pp
\item 正定性 $(a,a)\geq0$,且等号成立当且仅当$a=0$,
\end{enumerate}
\pp
则称$(a,b)$为$a$和$b$的\blue{内积}.
\pp
定义内积的复线性空间称为\blue{酉空间}.
\end{definition}
\pp
注: $(\lambda a,b)=\overline{\lambda}(a,b)$.
\end{frame}





\subsubsection{正定复矩阵*}
\begin{frame}\frametitle{正定复矩阵*}
\begin{definition}设$V$为$n$维酉空间.$G$为$n$阶复矩阵.
\pp
\begin{enumerate}
\item 称$G$为\blue{正定复矩阵},
\pp
若
\begin{enumerate}
\item $G^H=G$;
\pp
\item 任意给定列向量$z\in\bC^n$都有$z^HGz\geq0$,并且等号成立当且仅当$z=0$.
\end{enumerate}
\pp
\item 向量$a$的\blue{长度(或模长)}定义为 $|a|:=\sqrt{(a,a)}$.
\pp
\item \blue{垂直(正交)},\blue{标准正交基},\blue{Schmit正交化},\blue{正交补},\blue{Schwarz不等式},$\cdots$
\end{enumerate}
\end{definition}
\end{frame}





\subsubsection{共轭变换(伴随变换)*}
\begin{frame}\frametitle{共轭变换(伴随变换)*}
\pp
\begin{definition}[共轭变换,伴随变换]
设$\sA$为酉空间$V$上的一个线性变换.
\pp
则
\begin{itemize}
\item \blue{存在唯一}的$V$上的线性变换$\sA^*$满足对任意$a,b\in V$
\[(\sA a,b)=(a,\sA^*b).\]
\pp
称$\sA^*$是$\sA$的\blue{共轭变换}(或\blue{伴随变换}).
\pp
\item 若$\sA$在标准正交基下矩阵为$A$,则$\sA^*$在同一组标准正交基下矩阵为$A^H$.
\end{itemize}
\end{definition}
\end{frame}





\subsubsection{酉变换*}
\begin{frame}\frametitle{酉变换*}
设$\sU$为酉空间$V$上的线性变换.
\pp
则
\begin{equation*}
\begin{split}
\text{$\sU$为\blue{酉变换}} &\xLeftrightarrow{\text{定义}} \text{$\sU$保持内积}\\
\pp
&\Leftrightarrow \text{$\sU$保持向量长度}\\
\pp
&\Leftrightarrow \text{$\sU$将标准正交基变为标准正交基}\\
\pp
&\Leftrightarrow \sU^*\sU=\mathrm{id}_V\\
\pp
&\Leftrightarrow \text{$\sU$在标准正交基下的矩阵$U$为\blue{酉矩阵}(即,\blue{$U^HU=I$})}\\
\end{split}
\end{equation*}
\pp

\begin{theorem}
酉空间$V$上的全体酉变换组成的集合$U(V)$在复合的作用下构成群.
\end{theorem}
\end{frame}





\subsubsection{Hermite变换*}
\begin{frame}\frametitle{Hermite变换*}
设酉空间$V$上的线性变换$\sA$在某组标准正交基下矩阵为$A$.
\pp
则\footnote{在物理中,Hermite变换代表可观测量.}
\begin{equation*}
\begin{split}
\text{$\sA$为\blue{Hermite变换}} &\xLeftrightarrow{\text{定义}} \text{$\sA^*=\sA$(\blue{自伴随})}\\
\pp
&\Leftrightarrow (\sA a,b)=(a,\sA b),\quad (\forall a,b\in V)\\
\pp
&\Leftrightarrow \text{$A$为\blue{Hermite矩阵}(即, \blue{$A^H=A$}).}\\
\end{split}
\end{equation*}
\end{frame}





\subsubsection{规范变换*}
\begin{frame}\frametitle{规范变换*}
设$\sA$为酉空间$V$上的线性变换$\sA$.
\pp
\begin{equation*}
\begin{split}
\text{$\sA$为\blue{规范变换}} &\xLeftrightarrow{\text{定义}} \text{$\sA^*\sA=\sA\sA^*$}\\
\pp
&\Leftrightarrow A\text{酉相似于对角阵}\\
\pp
&\Leftrightarrow \text{$A$在某组标准正交基下为对角阵.}\\
\end{split}
\end{equation*}
\pp
\begin{corollary}设$\sU$为酉变换.
\pp
则
\begin{enumerate}
\item $\sU$的特征值$\lambda$模为$1$.即,存在$\theta\in[0,2\pi)$使得$\lambda=e^{i\theta}$.
\pp
\item $\sU$在某组标准正交基下矩阵为对角阵$\mathrm{diag}(e^{i\theta_1},e^{i\theta_2},\cdots,e^{i\theta_n})$.
\end{enumerate}
\end{corollary}
\pp
\yang{证明思路:$D^HD=I$ $\Rightarrow$ $|\lambda_i|=1$.}
\pp
\begin{corollary}设$\sA$为Hermite变换.
\pp
则
\begin{enumerate}
\item $\sA$的特征值全为实数;
\pp
\item $\sA$在某组标准正交基下的矩阵为实对角矩阵.
\end{enumerate}
\end{corollary}
\pp
\yang{证明思路:$D^H=D$ $\Rightarrow$ $\lambda_i\in\bR$.}

\pp
\blue{问题:} 反Hermite变换?
\end{frame}




